\documentclass[12pt]{article}
\usepackage[utf8]{inputenc}
\usepackage[T1]{fontenc}
\usepackage[brazil]{babel}
\usepackage[a4paper,margin=2.5cm]{geometry}
\usepackage{enumitem}
\setlist[itemize]{noitemsep,topsep=4pt}
\title{Sprint 3 - S3\_19\\Backend e Persistência}
\author{Matheus Moreli Durão Silva\\Matrícula 12621CBS217\\Cibersegurança - Bacharelado - Noturno}
\date{2026}
\begin{document}
\maketitle

\section*{Entregável}
MER, SQL e CRUD persistido. Este documento registra o objetivo do card, a solução desenvolvida e a forma de conferência. Os arquivos de código e as evidências citadas permanecem na mesma pasta deste documento.

\section*{O que foi estudado e aplicado}
\begin{itemize}
\item Entidades e relacionamentos
\item DDL SQL
\item SELECT, INSERT, UPDATE e DELETE
\item Validação de persistência
\end{itemize}

\section*{Como conferir}
O modelo 1:N está em \texttt{mer.md}; as tabelas e chaves estrangeiras estão em \texttt{schema.sql}. Execute \texttt{python testar\_api.py} para testar criação, leitura, alteração e exclusão com persistência SQLite. A API Flask em \texttt{api.py} habilita \texttt{PRAGMA foreign\_keys=ON} para aplicar a relação também nas requisições. O frontend de S3\_17 consome esta API.

\section*{Conclusão}
O artefato foi organizado para demonstrar tanto o conceito estudado quanto sua aplicação prática. As evidências de execução só são consideradas válidas quando produzidas pela ferramenta correspondente.
\end{document}
